\documentclass{amsart}
% For dealing with references we use the comment environment
\usepackage{verbatim}
\newenvironment{reference}{\comment}{\endcomment}
%\newenvironment{reference}{}{}
\newenvironment{slogan}{\comment}{\endcomment}
\newenvironment{history}{\comment}{\endcomment}

% For commutative diagrams we use Xy-pic
\usepackage[all]{xy}

% We use 2cell for 2-commutative diagrams.
\xyoption{2cell}
\UseAllTwocells

% We use multicol for the list of chapters between chapters
\usepackage{multicol}

% This is generally recommended for better output
\usepackage{lmodern}
\usepackage[T1]{fontenc}

% For cross-file-references

% Package for hypertext links:
\usepackage{hyperref}

% For any local file, say "hello.tex" you want to link to please

% Theorem environments.
%
\theoremstyle{plain}
\newtheorem{theorem}[subsection]{Theorem}
\newtheorem{proposition}[subsection]{Proposition}
\newtheorem{lemma}[subsection]{Lemma}

\theoremstyle{definition}
\newtheorem{definition}[subsection]{Definition}
\newtheorem{example}[subsection]{Example}
\newtheorem{exercise}[subsection]{Exercise}
\newtheorem{situation}[subsection]{Situation}

\theoremstyle{remark}
\newtheorem{remark}[subsection]{Remark}
\newtheorem{remarks}[subsection]{Remarks}

\numberwithin{equation}{subsection}

% Macros
%
\def\lim{\mathop{\mathrm{lim}}\nolimits}
\def\colim{\mathop{\mathrm{colim}}\nolimits}
\def\Spec{\mathop{\mathrm{Spec}}}
\def\Hom{\mathop{\mathrm{Hom}}\nolimits}
\def\Ext{\mathop{\mathrm{Ext}}\nolimits}
\def\SheafHom{\mathop{\mathcal{H}\!\mathit{om}}\nolimits}
\def\SheafExt{\mathop{\mathcal{E}\!\mathit{xt}}\nolimits}
\def\Sch{\mathit{Sch}}
\def\Mor{\mathop{\mathrm{Mor}}\nolimits}
\def\Ob{\mathop{\mathrm{Ob}}\nolimits}
\def\Sh{\mathop{\mathit{Sh}}\nolimits}
\def\NL{\mathop{N\!L}\nolimits}
\def\CH{\mathop{\mathrm{CH}}\nolimits}
\def\proetale{{pro\text{-}\acute{e}tale}}
\def\etale{{\acute{e}tale}}
\def\QCoh{\mathit{QCoh}}
\def\Ker{\mathop{\mathrm{Ker}}}
\def\Im{\mathop{\mathrm{Im}}}
\def\Coker{\mathop{\mathrm{Coker}}}
\def\Coim{\mathop{\mathrm{Coim}}}

% Boxtimes
%
\DeclareMathSymbol{\boxtimes}{\mathbin}{AMSa}{"02}

%
% Macros for moduli stacks/spaces
%
\def\QCohstack{\mathcal{QC}\!\mathit{oh}}
\def\Cohstack{\mathcal{C}\!\mathit{oh}}
\def\Spacesstack{\mathcal{S}\!\mathit{paces}}
\def\Quotfunctor{\mathrm{Quot}}
\def\Hilbfunctor{\mathrm{Hilb}}
\def\Curvesstack{\mathcal{C}\!\mathit{urves}}
\def\Polarizedstack{\mathcal{P}\!\mathit{olarized}}
\def\Complexesstack{\mathcal{C}\!\mathit{omplexes}}
% \Pic is the operator that assigns to X its picard group, usage \Pic(X)
% \Picardstack_{X/B} denotes the Picard stack of X over B
% \Picardfunctor_{X/B} denotes the Picard functor of X over B
\def\Pic{\mathop{\mathrm{Pic}}\nolimits}
\def\Picardstack{\mathcal{P}\!\mathit{ic}}
\def\Picardfunctor{\mathrm{Pic}}
\def\Deformationcategory{\mathcal{D}\!\mathit{ef}}

\setlength{\emergencystretch}{3em}
\begin{document}
\title[Stacks Project, Tag 0D2P]{631 --- Relative ampleness descends under fpqc coverings}
\maketitle
\noindent Copyright (C) 2005--2025 Johan de Jong. Stacks Project authors. Source: \href{https://stacks.math.columbia.edu/tag/0D2P}{Tag 0D2P}.

\noindent Editorial difficulty note (not author text). Difficulty H2: The proof reconstructs a relative Proj from the graded pushforward algebra, uses flat base change and faithful-flat stalk detection to establish the universal section map everywhere, then descends its open immersion. This coordinated descent of relative ampleness requires advanced scheme geometry..

\noindent Editorial provenance note (not author text). History: excerpt from revision \texttt{a0444\allowbreak 6e57e\allowbreak c1fbc\allowbreak 252a8\allowbreak 71afc\allowbreak ec775\allowbreak 2fb28\allowbreak 07b14}, descent.tex, lines 6532--6605. Reference presentation was adapted; original-author context and core support blocks were retained or added. The mathematical wording is unchanged. Licensed under GNU FDL 1.2 or later; no invariant sections or cover texts. See ../licenses/COPYING. Context blocks reproduce author definitions and supporting results; they are not additional counted problems.

\begin{definition}
\label{morphisms-definition-relatively-ample}
\begin{reference}
\href{https://stacks.math.columbia.edu/bibliography}{Bibliography: EGA (II Definition 4.6.1)}
\end{reference}
Let $f : X \to S$ be a morphism of schemes.
Let $\mathcal{L}$ be an invertible $\mathcal{O}_X$-module.
We say $\mathcal{L}$ is {\it relatively ample}, or {\it $f$-relatively ample},
or {\it ample on $X/S$}, or {\it $f$-ample} if $f : X \to S$
is quasi-compact, and if for every affine open $V \subset S$
the restriction of $\mathcal{L}$ to the open subscheme
$f^{-1}(V)$ of $X$ is ample.
\end{definition}

\begin{lemma}
\label{lemma-descending-property-open-immersion}
The property $\mathcal{P}(f) =$``$f$ is an open immersion''
is fpqc local on the base.
\end{lemma}

\begin{proof}
The property of being an open immersion is stable under base change,
see Schemes, Lemma \href{https://stacks.math.columbia.edu/tag/01JY}{lemma base change immersion (Tag 01JY)}.
The property of being an open immersion is Zariski local on the base
(this is obvious).

\medskip\noindent
Let $S' \to S$ be a flat surjective morphism of affine schemes,
and let $f : X \to S$ be a morphism. Assume that the base change
$f' : X' \to S'$ is an open immersion. We claim that $f$ is an
open immersion.
Then $f'$ is universally open, and universally injective.
Hence we conclude that $f$ is universally open by
Lemma \href{https://stacks.math.columbia.edu/tag/02KT}{lemma descending property universally open (Tag 02KT)}, and
universally injective by
Lemma \href{https://stacks.math.columbia.edu/tag/02KW}{lemma descending property universally injective (Tag 02KW)}.
In particular $f(X) \subset S$ is open. If for every affine
open $U \subset f(X)$ we can prove that $f^{-1}(U) \to U$
is an isomorphism, then $f$ is an open immersion and we're done.
If $U' \subset S'$ denotes the inverse image of $U$,
then $U' \to U$ is a faithfully flat morphism of affines and
$(f')^{-1}(U') \to U'$ is an isomorphism (as $f'(X')$ contains $U'$
by our choice of $U$). Thus we reduce to the case discussed
in the next paragraph.

\medskip\noindent
Let $S' \to S$ be a flat surjective morphism of affine schemes,
let $f : X \to S$ be a morphism, and assume that the base change
$f' : X' \to S'$ is an isomorphism. We have to show that $f$ is an
isomorphism also. It is clear that $f$ is surjective, universally injective,
and universally open (see arguments above for the last two).
Hence $f$ is bijective, i.e., $f$ is a homeomorphism.
Thus $f$ is affine by
Morphisms, Lemma \href{https://stacks.math.columbia.edu/tag/04DE}{lemma homeomorphism affine (Tag 04DE)}.
Since
$$
\mathcal{O}(S') \to
\mathcal{O}(X') =
\mathcal{O}(S') \otimes_{\mathcal{O}(S)} \mathcal{O}(X)
$$
is an isomorphism and since $\mathcal{O}(S) \to \mathcal{O}(S')$
is faithfully flat this implies that $\mathcal{O}(S) \to \mathcal{O}(X)$
is an isomorphism. Thus $f$ is an isomorphism. This finishes the proof of
the claim above.
Therefore Lemma \href{https://stacks.math.columbia.edu/tag/02KP}{lemma descending properties morphisms (Tag 02KP)} applies and we win.
\end{proof}


\begin{lemma}
\label{lemma-descending-property-ample}
Let $f : X \to S$ be a morphism of schemes.
Let $\mathcal{L}$ be an invertible $\mathcal{O}_X$-module.
Let $\{g_i : S_i \to S\}_{i \in I}$ be an fpqc covering.
Let $f_i : X_i \to S_i$ be the base change of $f$ and let $\mathcal{L}_i$
be the pullback of $\mathcal{L}$ to $X_i$.
The following are equivalent
\begin{enumerate}
\item $\mathcal{L}$ is ample on $X/S$, and
\item $\mathcal{L}_i$ is ample on $X_i/S_i$
for every $i \in I$.
\end{enumerate}
\end{lemma}

\begin{proof}
The implication (1) $\Rightarrow$ (2) follows from
Morphisms, Lemma \href{https://stacks.math.columbia.edu/tag/0893}{lemma ample base change (Tag 0893)}.
Assume $\mathcal{L}_i$ is ample on $X_i/S_i$ for every $i \in I$.
By Morphisms, Definition \ref{morphisms-definition-relatively-ample}
this implies that $X_i \to S_i$ is quasi-compact and by
Morphisms, Lemma \href{https://stacks.math.columbia.edu/tag/01VI}{lemma relatively ample separated (Tag 01VI)}
this implies $X_i \to S$ is separated.
Hence $f$ is quasi-compact and separated by
Lemmas \href{https://stacks.math.columbia.edu/tag/02KQ}{lemma descending property quasi compact (Tag 02KQ)} and
\href{https://stacks.math.columbia.edu/tag/02KU}{lemma descending property separated (Tag 02KU)}.

\medskip\noindent
This means that
$\mathcal{A} = \bigoplus_{d \geq 0} f_*\mathcal{L}^{\otimes d}$
is a quasi-coherent graded $\mathcal{O}_S$-algebra
(Schemes, Lemma \href{https://stacks.math.columbia.edu/tag/01LC}{lemma push forward quasi coherent (Tag 01LC)}).
Moreover, the formation of $\mathcal{A}$ commutes with flat
base change by
Cohomology of Schemes, Lemma \href{https://stacks.math.columbia.edu/tag/02KH}{lemma flat base change cohomology (Tag 02KH)}.
In particular, if we set
$\mathcal{A}_i = \bigoplus_{d \geq 0} f_{i, *}\mathcal{L}_i^{\otimes d}$
then we have $\mathcal{A}_i = g_i^*\mathcal{A}$.
It follows that the natural maps
$\psi_d : f^*\mathcal{A}_d \to \mathcal{L}^{\otimes d}$
of $\mathcal{O}_X$
pullback to give the natural maps
$\psi_{i, d} : f_i^*(\mathcal{A}_i)_d \to \mathcal{L}_i^{\otimes d}$
of $\mathcal{O}_{X_i}$-modules. Since $\mathcal{L}_i$ is ample on $X_i/S_i$
we see that for any point $x_i \in X_i$, there exists a $d \geq 1$
such that $f_i^*(\mathcal{A}_i)_d \to \mathcal{L}_i^{\otimes d}$
is surjective on stalks at $x_i$. This follows either directly
from the definition of a relatively ample module or from
Morphisms, Lemma \href{https://stacks.math.columbia.edu/tag/01VJ}{lemma characterize relatively ample (Tag 01VJ)}.
If $x \in X$, then we can choose an $i$ and an $x_i \in X_i$
mapping to $x$. Since $\mathcal{O}_{X, x} \to \mathcal{O}_{X_i, x_i}$
is flat hence faithfully flat, we conclude that for every $x \in X$
there exists a $d \geq 1$ such that
$f^*\mathcal{A}_d \to \mathcal{L}^{\otimes d}$
is surjective on stalks at $x$.
This implies that the open subset $U(\psi) \subset X$ of
Constructions, Lemma
\href{https://stacks.math.columbia.edu/tag/01O9}{lemma invertible map into relative proj (Tag 01O9)}
corresponding to the map
$\psi : f^*\mathcal{A} \to \bigoplus_{d \geq 0} \mathcal{L}^{\otimes d}$
of graded $\mathcal{O}_X$-algebras
is equal to $X$. Consider the corresponding morphism
$$
r_{\mathcal{L}, \psi} : X \longrightarrow \underline{\text{Proj}}_S(\mathcal{A})
$$
It is clear from the above that the base change of
$r_{\mathcal{L}, \psi}$ to $S_i$ is the morphism
$r_{\mathcal{L}_i, \psi_i}$ which is an open immersion by
Morphisms, Lemma \href{https://stacks.math.columbia.edu/tag/01VJ}{lemma characterize relatively ample (Tag 01VJ)}.
Hence $r_{\mathcal{L}, \psi}$ is an open immersion
by Lemma \ref{lemma-descending-property-open-immersion}
and we conclude $\mathcal{L}$ is ample on $X/S$ by
Morphisms, Lemma \href{https://stacks.math.columbia.edu/tag/01VJ}{lemma characterize relatively ample (Tag 01VJ)}.
\end{proof}
\end{document}
